%% toc: P15 157 157.3,157.4,157.5,157.6
\section{Lebesgue integral notation and the change of variables formula: proofs}
\label{sec:annex-p15-w02c-p15-001-integral-notation}

Use the notation of
Chapter~\ref{sec:p15-w02c-p15-001-integral-notation}, in particular
\cref{def:p15-w02c-p15-001-dx-notation}. Throughout, \(\lambda\) is Borel
Lebesgue measure on \(\mathbb R\). By
\cref{thm:p15-lebesgue-unit-interval}, every singleton is
\(\lambda\)-null and every bounded interval has its usual length.

\begin{annexitem}{annex:p15-w02c-p15-001-lebesgue-rectangles}{thm:p15-w02c-p15-001-lebesgue-rectangles}{Rectangle volumes determine Lebesgue measure}
\begin{proof}
Let \(\mathcal P\) consist of the empty set and all finite half-open
rectangles \(R=\prod_{i=1}^d(a_i,b_i]\) with real \(a_i<b_i\).

\paragraph{A pi-system.}
If \(R=\prod_i(a_i,b_i]\) and \(R'=\prod_i(a_i',b_i']\), then
\(R\cap R'=\prod_i(\max(a_i,a_i'),\min(b_i,b_i')]\). If
\(\max(a_i,a_i')\geq\min(b_i,b_i')\) for some \(i\), this set is empty;
otherwise it is again a finite half-open rectangle. Intersections with
\(\emptyset\) are empty. So \(\mathcal P\) is closed under finite
intersections.

\paragraph{It generates the Borel sets.}
An open rectangle \(\prod_i(c_i,e_i)\) is open: for a point \(x\) in it,
the Euclidean ball about \(x\) of radius
\(\min_i\min(x_i-c_i,e_i-x_i)>0\) lies in it, since each coordinate of a
point of the ball differs from \(x_i\) by less than that radius. Each
\(R\in\mathcal P\) is the countable intersection
\(R=\bigcap_{k\geq1}\prod_i(a_i,b_i+1/k)\) of open sets, hence Borel; so
\(\sigma(\mathcal P)\subseteq\mathcal B(\mathbb R^d)\). Conversely, for an
open rectangle with rational endpoints \(c_i<e_i\), choose \(k_0\) with
\(1/k_0<\min_i(e_i-c_i)\). Then
\(\prod_i(c_i,e_i)=\bigcup_{k\geq k_0}\prod_i(c_i,e_i-1/k]\), a countable
union of members of \(\mathcal P\). By \cref{thm:p15-borel-generators}
these open rectangles generate \(\mathcal B(\mathbb R^d)\), so
\(\mathcal B(\mathbb R^d)\subseteq\sigma(\mathcal P)\).

\paragraph{Agreement and exhaustion.}
By \cref{thm:p15-finite-product-measures} and the hypothesis, \(\lambda^d\)
and \(\nu\) both give each finite half-open rectangle its volume, and both
give \(\emptyset\) measure zero. The sets \(A_n=(-n,n]^d\), \(n\geq1\),
belong to \(\mathcal P\), increase, have union \(\mathbb R^d\), and have
measure \((2n)^d<\infty\) under both measures. Durrett's
uniqueness theorem for a pi-system with a finite exhaustion,
\cref{thm:p15-w02c-p15-002-sigma-finite-uniqueness}, therefore gives
\(\nu(B)=\lambda^d(B)\) for every \(B\in\sigma(\mathcal P)=\mathcal B(\mathbb R^d)\).
\end{proof}
\end{annexitem}

\begin{annexitem}{annex:p15-w02c-p15-001-stieltjes-notation}{thm:p15-w02c-p15-001-stieltjes-notation}{The measure named by \(dG(x)\)}
\begin{proof}
\paragraph{Part (1): monotonicity and right continuity.}
For real \(a<b\), \(G(b)-G(a)=\mu((a,b])\geq0\), so \(G\) is
nondecreasing. Fix \(b\in\mathbb R\) and put \(a=b-1\) and
\(A_n=(a,b+1/n]\) for \(n\geq1\). The sets \(A_n\) decrease, their
intersection is \((a,b]\), and \(\mu(A_1)=G(b+1)-G(a)\) is finite. By
part~(4) of \cref{thm:p15-measure-laws},
\(G(b+1/n)-G(a)=\mu(A_n)\to\mu((a,b])=G(b)-G(a)\), so
\(G(b+1/n)\to G(b)\). Because \(G\) is nondecreasing, for \(x\) with
\(b<x\leq b+1/n\) we have \(G(b)\leq G(x)\leq G(b+1/n)\); hence
\(G(x)\to G(b)\) as \(x\downarrow b\), and \(G\) is right-continuous.
Being real valued as well, \(G\) is a Stieltjes measure function in the
sense of \cref{def:p15-stieltjes-intervals}, and
\cref{thm:p15-stieltjes-construction} provides \(\mu_G\) with
\(\mu_G(I(a,b))=G(b)-G(a)\); for finite \(a<b\),
\(I(a,b)=(a,b]\).

\paragraph{Part (1): \(\mu=\mu_G\).}
Let \(\mathcal P\) consist of \(\varnothing\) and the intervals
\((a,b]\) with real \(a<b\). For two such intervals,
\((a,b]\cap(c,e]=(\max(a,c),\min(b,e)]\), which is empty when
\(\max(a,c)\geq\min(b,e)\) and otherwise lies in \(\mathcal P\); so
\(\mathcal P\) is a pi-system. Each \((a,b]=(-\infty,b]\setminus(-\infty,a]\)
is Borel, because the closed rays are among the generators of
\(\mathcal B(\mathbb R)\) listed in \cref{thm:p15-borel-generators}; so
\(\sigma(\mathcal P)\subseteq\mathcal B(\mathbb R)\). Conversely, for real
\(c\), \((-\infty,c]=\bigcup_{n\geq1}(c-n,c]\) lies in
\(\sigma(\mathcal P)\), and the same theorem says that these rays generate
\(\mathcal B(\mathbb R)\); so \(\mathcal B(\mathbb R)\subseteq\sigma(\mathcal P)\).
The measures \(\mu\) and \(\mu_G\) agree on \(\mathcal P\): both give
\(\varnothing\) measure zero and \((a,b]\) measure \(G(b)-G(a)\). The sets
\(A_n=(-n,n]\in\mathcal P\) increase to \(\mathbb R\), and both measures
give them the finite value \(G(n)-G(-n)\). By
\cref{thm:p15-w02c-p15-002-sigma-finite-uniqueness}, with both
sigma-algebras equal to \(\mathcal B(\mathbb R)\), \(\mu=\mu_G\).

\paragraph{Part (2).}
This is the defining property of \(\mu_G\) in
\cref{thm:p15-stieltjes-construction}, read with \(I(a,b)=(a,b]\) for
finite \(a<b\).

\paragraph{Part (3).}
The law \(\mathbb P_X\) is a measure on \(\mathcal B(\mathbb R)\)
(\cref{def:p15-pushforward-law}), and for real \(a<b\),
\(\mathbb P_X((a,b])=\mathbb P(a<X\leq b)=F_X(b)-F_X(a)\) by
\cref{thm:p15-cdf-laws}. Part~(1) with \(G=F_X\) and \(\mu=\mathbb P_X\)
gives \(\mathbb P_X=\mu_{F_X}\), and then
\(\int f(x)\,dF_X(x)=\int f\,d\mu_{F_X}=\int f\,d\mathbb P_X\) by
item~(5) of \cref{def:p15-w02c-p15-001-dx-notation}.

\paragraph{Part (4).}
By \cref{thm:p15-lebesgue-unit-interval}, \(\lambda\) is a Borel measure
with \(\lambda((a,b])=b-a\) for real \(a<b\). Part~(1) with
\(G(x)=x\) and \(\mu=\lambda\) gives \(\lambda=\mu_G\). Hence
\(\int f(x)\,dG(x)=\int f\,d\lambda=\int f(x)\,dx\), the last by item~(1)
of \cref{def:p15-w02c-p15-001-dx-notation} with \(d=1\) and
\(\lambda^1=\lambda\).

\paragraph{Part (5).}
The Dirac measure \(\delta_0\) of \cref{def:p15-probability-density} is a
Borel measure. Let \(a<b\) be real. If \(a<0\leq b\), then
\(0\in(a,b]\) and \(G(b)-G(a)=1-0=1=\delta_0((a,b])\). Otherwise
\(0\notin(a,b]\), so \(\delta_0((a,b])=0\), and either \(b<0\), when
\(G(a)=G(b)=0\), or \(a\geq0\), when \(G(a)=G(b)=1\); in both cases
\(G(b)-G(a)=0\). Part~(1) gives \(\mu_G=\delta_0\).

Now let \(f:\mathbb R\to[0,\infty]\) be Borel. Pointwise,
\(f=f\mathbf1_{\{0\}}+f\mathbf1_{\mathbb R\setminus\{0\}}\). The second
function is zero outside the Borel set \(\mathbb R\setminus\{0\}\), which
has \(\delta_0\)-measure zero, so it equals zero \(\delta_0\)-a.e. and its
integral is zero by part~(2) of \cref{thm:p15-nonnegative-integral-laws}.
Part~(1) of \cref{thm:p15-w02c-p15-001-counting-sums}, with \(s=0\),
gives \(\int f\mathbf1_{\{0\}}\,d\delta_0=f(0)\,\delta_0(\{0\})=f(0)\).
Additivity of the nonnegative integral
(\cref{thm:p15-nonnegative-integral-laws}) gives
\(\int f\,d\delta_0=f(0)\). For a real Borel \(f\), apply this to
\(|f|\), \(f^+\) and \(f^-\): \(\int|f|\,d\delta_0=|f(0)|<\infty\), so
\(f\) is integrable, and
\(\int f\,d\delta_0=f^+(0)-f^-(0)=f(0)\) by
\cref{def:p15-integrable-functions}.
\end{proof}
\end{annexitem}

\begin{annexitem}{annex:p15-w02c-p15-001-counting-sums}{thm:p15-w02c-p15-001-counting-sums}{Counting measure}
\begin{proof}
\paragraph{Part (1).}
Put \(c=f(s)\in[0,\infty]\). Pointwise, \(f\mathbf1_{\{s\}}=c\mathbf1_{\{s\}}\):
both are zero off \(s\) and equal \(c\) at \(s\), with the convention
of \cref{def:p15-nonnegative-integral} for \(f\mathbf1_A\). If
\(c<\infty\), this is the nonnegative simple function with parts
\(\{s\}\) and \(\Omega\setminus\{s\}\) and coefficients \(c\) and \(0\);
by \cref{def:p15-simple-integral} and the agreement of the two integrals
in \cref{thm:p15-monotone-integral}, its integral is
\(c\,\mu(\{s\})+0\cdot\mu(\Omega\setminus\{s\})=c\,\mu(\{s\})\). If
\(c=\infty\), the simple functions \(n\mathbf1_{\{s\}}\) increase
pointwise to \(\infty\mathbf1_{\{s\}}\), so by monotone convergence the
integral is \(\lim_nn\,\mu(\{s\})\). This limit is \(0\) when
\(\mu(\{s\})=0\) and \(\infty\) when \(\mu(\{s\})>0\); in both cases it
equals \(\infty\cdot\mu(\{s\})\) under the products of
\cref{def:p15-simple-integral}.

\paragraph{Part (2).}
Every function on \(I\) is measurable, since every subset of \(I\) is.
Let \((i_k)\) be an enumeration of \(I\). The functions
\(f\mathbf1_{\{i_k\}}\) are nonnegative and measurable, and
\(f=\sum_kf\mathbf1_{\{i_k\}}\) pointwise, because each \(i\in I\) equals
exactly one \(i_k\), and every other term vanishes at \(i\). When \(I\) is
infinite, the countable-sum identity of
\cref{thm:p15-nonnegative-integral-laws} gives
\(\int f\,d\#=\sum_k\int f\mathbf1_{\{i_k\}}\,d\#\); when \(I\) is finite
with \(m\) elements, finite additivity in the same theorem, applied
\(m-1\) times, gives the same identity with a finite sum. By part~(1)
and \(\#(\{i_k\})=1\), the \(k\)-th term is \(f(i_k)\). So
\(\int f\,d\#=\sum_kf(i_k)\), the limit in \([0,\infty]\) of the
nondecreasing partial sums \(\sum_{k\leq N}f(i_k)\) (for finite \(I\),
these are constant from \(N=m\) on).

Each partial sum is the sum of \(f\) over the finite set
\(\{i_1,\ldots,i_N\}\), so the limit is at most the supremum \(\sigma\) of
finite-subset sums. Conversely, a finite \(F\subseteq I\) is contained in
\(\{i_1,\ldots,i_N\}\) for \(N\) at least the largest index of a member of
\(F\), and since the terms are nonnegative,
\(\sum_{i\in F}f(i)\leq\sum_{k\leq N}f(i_k)\leq\int f\,d\#\). Taking the
supremum over \(F\) gives \(\sigma\leq\int f\,d\#\). Hence
\(\int f\,d\#=\sigma\), a quantity that does not mention an
enumeration.

\paragraph{Part (3).}
The functions \(|f|\), \(f^+\), \(f^-\) are nonnegative and real valued.
By part~(2), \(\int|f|\,d\#=\sum_{i\in I}|f(i)|\), so \(f\) is integrable
exactly when this sum is finite. Suppose it is, and let \((i_k)\) be an
enumeration. By part~(2), \(\sum_k|f(i_k)|=\int|f|\,d\#<\infty\), so the
series converges absolutely. Also by part~(2),
\(\sum_{k\leq N}f^\pm(i_k)\to\int f^\pm\,d\#\), and both limits are at most
\(\int|f|\,d\#<\infty\). Since \(f=f^+-f^-\) pointwise,
\[
 \sum_{k\leq N}f(i_k)=\sum_{k\leq N}f^+(i_k)-\sum_{k\leq N}f^-(i_k)
 \ \to\ \int f^+\,d\#-\int f^-\,d\#=\int f\,d\#,
\]
the last equality by \cref{def:p15-integrable-functions}.

\paragraph{Part (4).}
By \cref{thm:p15-monotone-integral} with \(\mu=\#\),
\(\int f_n\,d\#\uparrow\int f\,d\#\). By part~(2) these integrals are the
sums \(\sum_{i\in I}f_n(i)\) and \(\sum_{i\in I}f(i)\).
\end{proof}
\end{annexitem}

\begin{annexitem}{annex:p15-w02c-p15-001-interval-integrals}{thm:p15-w02c-p15-001-interval-integrals}{Endpoints and additivity of interval integrals}
\begin{proof}
\paragraph{Part (1).}
Let \(E\) be one of \((a,b)\), \((a,b]\) and \([a,b)\). Then
\(E\subseteq[a,b]\) and \([a,b]\setminus E\subseteq\{a,b\}\). The set
\(\{a,b\}\) is Borel and, by subadditivity
(\cref{thm:p15-measure-laws}),
\(\lambda(\{a,b\})\leq\lambda(\{a\})+\lambda(\{b\})=0\). Hence
\(f\mathbf1_{[a,b]}=f\mathbf1_E\) outside a \(\lambda\)-null Borel set.

Suppose \(f\geq0\) on \([a,b]\). Then \(f\mathbf1_{[a,b]}\) and
\(f\mathbf1_E\) are nonnegative Borel functions that agree a.e., so their
integrals agree by part~(2) of \cref{thm:p15-nonnegative-integral-laws}.
That is the claim.

Suppose instead that \(\int_a^b|f(x)|\,dx<\infty\). The positive parts
\(f^+\mathbf1_{[a,b]}\) and \(f^+\mathbf1_E\) are nonnegative and agree
a.e., so they have the same integral, and this integral is at most
\(\int|f|\mathbf1_{[a,b]}\,d\lambda<\infty\) by monotonicity
(\cref{thm:p15-monotone-integral}), since \(f^+\leq|f|\). The same holds
for the negative parts. Because \((f\mathbf1_E)^\pm=f^\pm\mathbf1_E\),
\cref{def:p15-integrable-functions} gives
\[
 \begin{aligned}
 \int_Ef\,d\lambda
 &=\int f^+\mathbf1_E\,d\lambda-\int f^-\mathbf1_E\,d\lambda\\
 &=\int f^+\mathbf1_{[a,b]}\,d\lambda-\int f^-\mathbf1_{[a,b]}\,d\lambda\\
 &=\int_a^bf(x)\,dx,
 \end{aligned}
\]
a difference of finite numbers. Also
\(\int|f|\mathbf1_E\,d\lambda\leq\int|f|\mathbf1_{[a,b]}\,d\lambda<\infty\),
so all four integrals are finite.

\paragraph{Part (2).}
Let \(a\leq c\leq b\). The sets \([a,c)\) and \([c,b]\) are disjoint
Borel sets with union \([a,b]\); the first is empty when \(c=a\).

Suppose \(f\geq0\) on \([a,b]\). Apply part~(4) of
\cref{thm:p15-nonnegative-integral-laws} to the nonnegative function
\(f\mathbf1_{[a,b]}\): the set function
\(A\mapsto\int_Af\mathbf1_{[a,b]}\,d\lambda\) is a measure, so
\[
 \int_a^bf(x)\,dx
 =\int_{[a,c)}f\,d\lambda+\int_{[c,b]}f\,d\lambda .
\]
The second term is \(\int_c^bf(x)\,dx\) by definition. The first equals
\(\int_a^cf(x)\,dx\) by part~(1) applied on \([a,c]\), where \(f\geq0\);
when \(c=a\) both are zero, since \([a,a)=\emptyset\) and part~(1) gives
\(\int_a^af(x)\,dx=\int_{[a,a)}f\,d\lambda=0\).

Suppose instead that \(\int_a^b|f(x)|\,dx<\infty\). Apply the nonnegative
case to \(f^+\) and to \(f^-\), which are nonnegative everywhere. Every
term is at most \(\int_a^b|f(x)|\,dx\) by monotonicity, so every term is
finite. Subtracting the identity for \(f^-\) from the identity for
\(f^+\) and using \cref{def:p15-integrable-functions} on each of the three
intervals gives the claim between finite numbers. The hypothesis of
part~(1) holds on \([a,c]\) and \([c,b]\) because
\(\int_a^c|f(x)|\,dx\) and \(\int_c^b|f(x)|\,dx\) are at most
\(\int_a^b|f(x)|\,dx\).

\paragraph{Part (3).}
If \(k\geq0\), then \(k\mathbf1_{[a,b]}\) is a nonnegative simple
function and \cref{thm:p15-simple-integral-laws} gives
\(\int k\mathbf1_{[a,b]}\,d\lambda=k\lambda([a,b])=k(b-a)\); for
\(a=b\) the interval is a singleton of measure zero, and the formula
still reads \(k(b-a)=0\). If \(k<0\), the function
\(k\mathbf1_{[a,b]}\) has positive part zero and negative part
\(|k|\mathbf1_{[a,b]}\), whose integral is \(|k|(b-a)<\infty\). Hence
its signed integral is \(0-|k|(b-a)=k(b-a)\).
\end{proof}
\end{annexitem}

\begin{annexitem}{annex:p15-w02c-p15-001-interval-limits}{thm:p15-w02c-p15-001-interval-limits}{Limits of interval integrals}
\begin{proof}
\paragraph{Part (1).}
If \(0\leq b\leq b'\), then \([0,b]\subseteq[0,b']\), so
\(f\mathbf1_{[0,b]}\leq f\mathbf1_{[0,b']}\) and
\(\int_0^bf(x)\,dx\leq\int_0^{b'}f(x)\,dx\) by monotonicity
(\cref{thm:p15-monotone-integral}). A nondecreasing function on
\([0,\infty)\) with values in \([0,\infty]\) has a limit in
\([0,\infty]\) as \(b\to\infty\), namely its supremum, and that limit is
also the limit along the integers \(b=k\). For each real \(x\),
\(f(x)\mathbf1_{[0,k]}(x)\) is zero for all \(k\) if \(x<0\), and equals
\(f(x)\) for all \(k\geq x\) if \(x\geq0\). Thus
\(f\mathbf1_{[0,k]}\uparrow f\mathbf1_{[0,\infty)}\) pointwise, and
monotone convergence gives
\(\int_0^kf(x)\,dx\uparrow\int_{[0,\infty)}f\,d\lambda\).

\paragraph{Part (2).}
For each real \(b\geq0\),
\(|f\mathbf1_{[0,b]}|\leq|f|\mathbf1_{[0,\infty)}\), whose integral is
finite; hence each \(\int_0^bf(x)\,dx\) is finite. Let \(b_k\) be any
sequence of nonnegative reals with \(b_k\to\infty\). For each real
\(x\), \(f(x)\mathbf1_{[0,b_k]}(x)\to f(x)\mathbf1_{[0,\infty)}(x)\): if
\(x<0\) every term is zero, and if \(x\geq0\) then \(b_k\geq x\) for all
large \(k\). The functions are dominated by the integrable
\(|f|\mathbf1_{[0,\infty)}\), so part~(1) of
\cref{thm:p15-dominated-bounded-convergence} gives
\(\int_0^{b_k}f(x)\,dx\to\int_{[0,\infty)}f\,d\lambda\). Since every
sequence tending to infinity gives the same limit, the limit as
\(b\to\infty\) exists and equals \(\int_{[0,\infty)}f\,d\lambda\).

\paragraph{Part (3): measurability and the interval formula.}
No \(c_n\) is zero. For a Borel set \(B\subseteq\mathbb R\),
\[
 g^{-1}(B)=\bigcup_{n\geq0:\ c_n\in B}[n,n+1)
 \ \cup\ \begin{cases}(-\infty,0),&0\in B,\\\emptyset,&0\notin B,\end{cases}
\]
a countable union of intervals, hence Borel. Also \(|g|\leq1\).

Fix real \(b\geq0\) and let \(N=\lfloor b\rfloor\), so
\(N\leq b<N+1\). The interval \([0,b]\) is the disjoint union of the
intervals \([n,n+1)\), \(0\leq n\leq N-1\), and \([N,b]\). On \([N,b]\)
the function \(g\) equals \(c_N\). Therefore
\[
 g\mathbf1_{[0,b]}
 =\sum_{n=0}^{N-1}c_n\mathbf1_{[n,n+1)}+c_N\mathbf1_{[N,b]},
\]
a finite linear combination of indicators of sets of finite measure. Each
indicator is integrable, so signed linearity
(\cref{thm:p15-signed-integral-laws}) and
\(\int\mathbf1_A\,d\lambda=\lambda(A)\) give
\[
 \int_0^bg(x)\,dx
 =\sum_{n=0}^{N-1}c_n\lambda([n,n+1))+c_N\lambda([N,b])
 =S_N+(b-N)c_N .
\]

\paragraph{Part (3): the limit.}
From \(S_{N+1}=S_N+c_N\), for every \(m\geq0\),
\[
 \begin{aligned}
 S_{2m+2}-S_{2m}&=\frac1{2m+1}-\frac1{2m+2}>0,\\
 S_{2m+3}-S_{2m+1}&=-\frac1{2m+2}+\frac1{2m+3}<0,\\
 S_{2m+1}-S_{2m}&=\frac1{2m+1}>0 .
 \end{aligned}
\]
Also \(S_2=1-\tfrac12=\tfrac12\) and
\(S_3=1-\tfrac12+\tfrac13=\tfrac56\). Hence, for every \(m\geq1\),
\(\tfrac12=S_2\leq S_{2m}<S_{2m+1}\leq S_3=\tfrac56\). The sequence
\((S_{2m})_{m\geq1}\) is nondecreasing and bounded above by \(5/6\), so
it converges to some \(L_e\geq1/2\); the sequence
\((S_{2m+1})_{m\geq1}\) is nonincreasing and bounded below by \(1/2\), so
it converges to some \(L_o\leq5/6\). Their difference
\(S_{2m+1}-S_{2m}=1/(2m+1)\) tends to zero, so \(L_e=L_o\). Call this
common value \(L\); then \(1/2\leq L\leq5/6\), and \(S_N\to L\) because
both the even-indexed and the odd-indexed terms converge to \(L\).

Let \(\varepsilon>0\) and choose \(N_0\) with \(|S_N-L|<\varepsilon\)
for all \(N\geq N_0\). For real \(b\geq N_0\), \(N=\lfloor b\rfloor\geq N_0\),
and with \(t=b-N\in[0,1)\),
\[
 \int_0^bg(x)\,dx=S_N+tc_N=(1-t)S_N+tS_{N+1},
\]
so
\(\bigl|\int_0^bg(x)\,dx-L\bigr|\leq(1-t)|S_N-L|+t|S_{N+1}-L|<\varepsilon\).
Thus \(\int_0^bg(x)\,dx\to L\) as \(b\to\infty\).

\paragraph{Part (3): the half-line integral is undefined.}
The positive and negative parts of \(g\) are
\[
 g^+=\sum_{k\geq0}\frac1{2k+1}\mathbf1_{[2k,2k+1)},\qquad
 g^-=\sum_{k\geq0}\frac1{2k+2}\mathbf1_{[2k+1,2k+2)} ,
\]
since \(c_n>0\) exactly for even \(n\). Both vanish on \((-\infty,0)\), so
\((g\mathbf1_{[0,\infty)})^\pm=g^\pm\). The countable-sum formula of
\cref{thm:p15-nonnegative-integral-laws}, with
\(\lambda([n,n+1))=1\), gives
\[
 \int_{[0,\infty)}g^+\,d\lambda=\sum_{k\geq0}\frac1{2k+1},\qquad
 \int_{[0,\infty)}g^-\,d\lambda=\sum_{k\geq0}\frac1{2k+2}
 =\frac12\sum_{j\geq1}\frac1j .
\]
The harmonic series diverges: for \(J\geq1\), grouping the terms with
\(2^i<j\leq2^{i+1}\), each of the \(2^i\) terms in a group is at least
\(2^{-(i+1)}\), so
\[
 \sum_{j=1}^{2^J}\frac1j
 =1+\sum_{i=0}^{J-1}\ \sum_{j=2^i+1}^{2^{i+1}}\frac1j
 \geq1+\sum_{i=0}^{J-1}2^i\,2^{-(i+1)}=1+\frac J2 .
\]
Hence \(\int_{[0,\infty)}g^-\,d\lambda=\infty\). Since
\(1/(2k+1)>1/(2k+2)\) for every \(k\), also
\(\int_{[0,\infty)}g^+\,d\lambda=\infty\). Both part integrals are infinite,
so by \cref{def:p15-integrable-functions} the signed integral
\(\int_{[0,\infty)}g\,d\lambda\) is undefined.
\end{proof}
\end{annexitem}

\begin{annexitem}{annex:p15-w02c-p15-001-change-of-variables}{thm:p15-w02c-p15-001-change-of-variables}{Change of variables formula}
\begin{proof}
For \(A\in\mathcal S\), the hypothesis and \cref{def:p15-pushforward-law}
give \(\mu(A)=\mathbb P(X\in A)=\mathbb P(X^{-1}(A))=\mathbb P_X(A)\).
With \(h=X\), the definition
\((\mathbb P\circ h^{-1})(A)=\mathbb P(h^{-1}(A))\) gives the same
number. Hence \(\mu=\mathbb P_X=\mathbb P\circ h^{-1}\) as measures on
\((S,\mathcal S)\).

The composite \(f(X)\) is a real random variable by
\cref{thm:p15-measurable-operations}. Apply
\cref{thm:p15-expectation-law-integral} with its measurable function taken
to be \(f\). If \(f\geq0\), or if \(\mathbb E|f(X)|<\infty\), that
theorem gives \(\mathbb E[f(X)]=\int_Sf\,d\mathbb P_X\), which is
\(\int_Sf(y)\,\mu(dy)\) in the notation of item~(4) of
\cref{def:p15-w02c-p15-001-dx-notation}. The same theorem gives
\(\mathbb E|f(X)|=\int_S|f|\,d\mathbb P_X\), so the two integrability
conditions are equivalent. Finally,
\[
 \mathbb E[f(X)]=\int_\Omega f(X(\omega))\,d\mathbb P(\omega)
\]
by \cref{def:p15-expectation-integral}; writing \(h\) for \(X\) and
\(\mathbb P\circ h^{-1}\) for \(\mu\) turns the identity just proved into
the second displayed form.

\paragraph{The same identity from the four-step theorem.}
For measurable \(h:S\to[0,\infty]\), put
\(\Phi(h)=\mathbb E[h(X)]=\int_\Omega h\circ X\,d\mathbb P\) and
\(\Psi(h)=\int_Sh\,d\mu\); the composite \(h\circ X\) is measurable by
\cref{thm:p15-measurable-operations}. Hypothesis~(a) of
\cref{thm:p15-w02c-p15-001-four-step} holds for both: pointwise
\((h+g)\circ X=h\circ X+g\circ X\) and \((ch)\circ X=c(h\circ X)\), and
the nonnegative integral on either space is additive and commutes with
finite constants \(c\geq0\) (\cref{thm:p15-nonnegative-integral-laws}).
Hypothesis~(b) holds for both: if \(h_n\uparrow h\) pointwise on \(S\),
then \(h_n\circ X\uparrow h\circ X\) pointwise on \(\Omega\), and
\cref{thm:p15-monotone-integral} applies on each space. For
\(B\in\mathcal S\), \(\mathbf1_B\circ X=\mathbf1_{X^{-1}(B)}\), so
\(\Phi(\mathbf1_B)=\mathbb P(X\in B)=\mu(B)=\Psi(\mathbf1_B)\) by
\cref{thm:p15-simple-integral-laws}. Part~(1) of the four-step theorem
gives \(\mathbb E f(X)=\int_Sf\,d\mu\) for \(f\geq0\). For real \(f\),
part~(2) gives the equivalence of \(\mathbb E|f(X)|<\infty\) and
\(\int_S|f|\,d\mu<\infty\), and, since \((f\circ X)^\pm=f^\pm\circ X\),
\(\mathbb E f(X)=\Phi(f^+)-\Phi(f^-)=\Psi(f^+)-\Psi(f^-)=\int_Sf\,d\mu\).
\end{proof}
\end{annexitem}

\begin{annexitem}{annex:p15-w02c-p15-001-uniform-integrals}{thm:p15-w02c-p15-001-uniform-integrals}{On the uniform variable}
\begin{proof}
\paragraph{The law of \(U\).}
For Borel \(B\), \(U^{-1}(B)=\{u\in(0,1):u\in B\}=B\cap(0,1)\), which
belongs to \(\mathcal F_U\). Hence
\(\mathbb U_U(B)=\mathbb U(B\cap(0,1))=\lambda(B\cap(0,1))\). Since
\(\mathbf1_B\mathbf1_{(0,1)}=\mathbf1_{B\cap(0,1)}\),
\cref{thm:p15-simple-integral-laws} gives
\(\int_B\mathbf1_{(0,1)}\,d\lambda=\lambda(B\cap(0,1))\). Taking
\(B=\mathbb R\) gives \(\int\mathbf1_{(0,1)}\,d\lambda=1\), so the
nonnegative Borel function \(\mathbf1_{(0,1)}\) is a probability density
for \(\mathbb U_U\) in the sense of \cref{def:p15-probability-density}.

\paragraph{Reduction to the open interval.}
Put \(\tilde h\defeq h\mathbf1_{(0,1)}\). It is Borel, since
\(\tilde h^{-1}(B)=(h^{-1}(B)\cap(0,1))\cup((0,1)^c)\) when \(0\in B\)
and \(\tilde h^{-1}(B)=h^{-1}(B)\cap(0,1)\) otherwise. Because
\(U(u)=u\in(0,1)\) for every outcome, \(\tilde h(U)=h(U)\) everywhere
on \(\Omega_U\).

\paragraph{The nonnegative case.}
Suppose \(h\geq0\) on \([0,1]\). Then \(\tilde h\geq0\). The change of
variables formula (\cref{thm:p15-w02c-p15-001-change-of-variables}) and
the density formula of \cref{thm:p15-density-integrals} give
\[
 \mathbb E[h(U)]=\mathbb E[\tilde h(U)]
 =\int\tilde h\,d\mathbb U_U
 =\int\tilde h\mathbf1_{(0,1)}\,d\lambda
 =\int_{(0,1)}h\,d\lambda,
\]
and part~(1) of \cref{thm:p15-w02c-p15-001-interval-integrals}, whose
hypothesis \(h\geq0\) on \([0,1]\) holds, identifies the last integral
with \(\int_0^1h(x)\,dx\). All terms lie in \([0,\infty]\). Applied to the
nonnegative function \(|h|\), this proves
\(\mathbb E[|h(U)|]=\int_0^1|h(x)|\,dx\) for every Borel \(h\).

\paragraph{The integrable case.}
Suppose \(\int_0^1|h(x)|\,dx<\infty\). By the identity just proved,
\(\mathbb E|\tilde h(U)|=\mathbb E|h(U)|<\infty\). The change of
variables formula therefore applies to \(\tilde h\), and its integrability
statement gives \(\int|\tilde h|\,d\mathbb U_U<\infty\). Apply the density
formula to the nonnegative functions \(\tilde h^+\) and \(\tilde h^-\):
\[
 \int\tilde h^\pm\,d\mathbb U_U=\int\tilde h^\pm\mathbf1_{(0,1)}\,d\lambda
 =\int h^\pm\mathbf1_{(0,1)}\,d\lambda ,
\]
both finite. Subtracting, and using
\((h\mathbf1_{(0,1)})^\pm=h^\pm\mathbf1_{(0,1)}\) with
\cref{def:p15-integrable-functions},
\[
 \mathbb E[h(U)]=\int\tilde h\,d\mathbb U_U=\int_{(0,1)}h\,d\lambda
 =\int_0^1h(x)\,dx,
\]
the last step by part~(1) of
\cref{thm:p15-w02c-p15-001-interval-integrals} in its integrable case.
\end{proof}
\end{annexitem}

\begin{annexitem}{annex:p15-w02c-p15-001-four-step}{thm:p15-w02c-p15-001-four-step}{The four-step method}
\begin{proof}
Let \(\Theta\) be either of \(\Phi\), \(\Psi\).

\paragraph{Step 1: indicators.}
This is the hypothesis \(\Phi(\mathbf1_B)=\Psi(\mathbf1_B)\) for
\(B\in\mathcal S\).

\paragraph{Step 2: simple functions.}
Let \(s=\sum_{j=1}^ma_j\mathbf1_{A_j}\) be a nonnegative simple function
as in \cref{def:p15-simple-integral}, with finite \(a_j\geq0\) and
\(A_j\in\mathcal S\). Induction on \(m\) using~(a) gives
\(\Theta(s)=\sum_{j=1}^ma_j\Theta(\mathbf1_{A_j})\): for \(m=1\) this is
the homogeneity in~(a), and
\(\Theta\bigl(\sum_{j\leq m}a_j\mathbf1_{A_j}\bigr)
=\Theta\bigl(\sum_{j<m}a_j\mathbf1_{A_j}\bigr)+a_m\Theta(\mathbf1_{A_m})\)
by additivity and homogeneity. By step~1 the two expansions agree term by
term, so \(\Phi(s)=\Psi(s)\).

\paragraph{Step 3: nonnegative functions.}
Let \(f\in\mathcal M_+\). If \(S=\varnothing\), then \(f=\mathbf1_\varnothing\)
and step~1 applies. Otherwise fix \(s_0\in S\) and let
\(\delta(B)=1\) if \(s_0\in B\) and \(\delta(B)=0\) otherwise, for
\(B\in\mathcal S\). This is a probability measure: \(\delta(S)=1\), and
of countably many disjoint sets at most one contains \(s_0\), so
\(\delta\) is countably additive. On the probability space
\((S,\mathcal S,\delta)\), \(f\) is a nonnegative extended real random
variable, so \cref{thm:p15-measurable-limits} provides nonnegative simple
functions \(s_n\) on \((S,\mathcal S)\) with \(s_n\leq s_{n+1}\) and
\(s_n\to f\) pointwise. Being simple and increasing to \(f\) are
properties of the functions and of \(\mathcal S\) alone; the measure
\(\delta\) plays no further part. By~(b) and step~2,
\[
 \Phi(f)=\lim_n\Phi(s_n)=\lim_n\Psi(s_n)=\Psi(f).
\]
This proves part~(1).

\paragraph{Step 4: real functions.}
Let \(f\) be real and measurable. The functions \(|f|\), \(f^+\), \(f^-\)
belong to \(\mathcal M_+\) (\cref{def:p15-integrable-functions}). By
part~(1), \(\Phi(|f|)=\Psi(|f|)\), so one is finite exactly when the
other is. Since \(|f|=f^++f^-\), (a) gives
\(\Theta(|f|)=\Theta(f^+)+\Theta(f^-)\), so when \(\Theta(|f|)<\infty\) both
\(\Theta(f^+)\) and \(\Theta(f^-)\) are finite. By part~(1),
\(\Phi(f^\pm)=\Psi(f^\pm)\), and subtracting the finite values gives
\(\Phi(f^+)-\Phi(f^-)=\Psi(f^+)-\Psi(f^-)\). This proves part~(2).
\end{proof}
\end{annexitem}

\begin{annexitem}{annex:p15-w02c-p15-001-cauchy-schwarz}{thm:p15-w02c-p15-001-cauchy-schwarz}{Discriminant proof}
\begin{proof}
\paragraph{Integrability of \(fg\).}
The functions \(f^2\), \(g^2\), \(fg\) and \(|fg|\) are measurable by
\cref{thm:p15-measurable-operations}. Pointwise,
\(0\leq(|f|-|g|)^2=f^2-2|fg|+g^2\), so \(|fg|\leq(f^2+g^2)/2\). By the
additivity and scaling laws of \cref{thm:p15-nonnegative-integral-laws}
and its order statement,
\[
 \int|fg|\,d\mu\leq\frac12\int f^2\,d\mu+\frac12\int g^2\,d\mu<\infty .
\]
So \(fg\) is integrable, its signed integral is a real number, and
\(a\), \(b\) and \(c\) are real.

\paragraph{Part (1).}
Fix real \(\theta\). Pointwise,
\((f+\theta g)^2=f^2+2\theta\,fg+\theta^2g^2\). The three functions
\(f^2\), \(fg\) and \(g^2\) are real, measurable and integrable, so
the linearity of \cref{thm:p15-signed-integral-laws} gives
\[
 \int(f+\theta g)^2\,d\mu
 =\int f^2\,d\mu+2\theta\int fg\,d\mu+\theta^2\int g^2\,d\mu
 =a\theta^2+b\theta+c .
\]
For the nonnegative functions \(f^2\) and \(g^2\), the signed integral
equals the nonnegative integral, since their negative parts are zero
(\cref{def:p15-integrable-functions}). The function
\((f+\theta g)^2\) is nonnegative, so for the same reason its integral is a
nonnegative integral and lies in \([0,\infty)\). Write
\(Q(\theta)\defeq a\theta^2+b\theta+c\); then \(Q(\theta)\geq0\) for every
real \(\theta\).

\paragraph{Part (2) when \(a>0\).}
Suppose, for contradiction, that \(b^2-4ac>0\), and let
\(D\defeq\sqrt{b^2-4ac}>0\). Put
\(\theta_\pm\defeq(-b\pm D)/(2a)\), the two distinct real roots given
by the formula for the roots of a quadratic. Then
\(\theta_++\theta_-=-b/a\) and
\(\theta_+\theta_-=(b^2-D^2)/(4a^2)=c/a\), so for every real \(\theta\),
\[
 a(\theta-\theta_+)(\theta-\theta_-)
 =a\theta^2-a(\theta_++\theta_-)\theta+a\theta_+\theta_-
 =a\theta^2+b\theta+c=Q(\theta).
\]
At the midpoint \(\theta_0\defeq-b/(2a)\) of the two roots,
\(\theta_0-\theta_+=-D/(2a)\) and \(\theta_0-\theta_-=D/(2a)\), so
\(Q(\theta_0)=-D^2/(4a)<0\). This contradicts part~(1). Hence
\(b^2-4ac\leq0\). This is the step in Durrett's remark from nonnegativity to at most one
real root: with \(a>0\), two distinct real roots would make the
quadratic negative between them.

\paragraph{Part (2) when \(a=0\).}
Then \(Q(\theta)=b\theta+c\geq0\) for every real \(\theta\). If
\(b\neq0\), the choice \(\theta=-(c+1)/b\) gives \(Q(\theta)=-1<0\),
which is impossible. Hence \(b=0\) and \(b^2-4ac=0\).

In both cases \(b^2\leq4ac\). Substituting the definitions,
\(4\left(\int fg\,d\mu\right)^2\leq4\int f^2\,d\mu\int g^2\,d\mu\), which
is the displayed inequality.

\paragraph{Part (3).}
The functions \(|f|\) and \(|g|\) are real and measurable, with
\(|f|^2=f^2\) and \(|g|^2=g^2\), so they satisfy the hypotheses. Part~(2)
for them, with \(|f||g|=|fg|\), gives
\(\left(\int|fg|\,d\mu\right)^2\leq\int f^2\,d\mu\int g^2\,d\mu\). Both
sides are nonnegative real numbers, and the square root is increasing on
\([0,\infty)\), so
\[
 \int|fg|\,d\mu\leq\left(\int f^2\,d\mu\right)^{1/2}
 \left(\int g^2\,d\mu\right)^{1/2}=\|f\|_2\|g\|_2
\]
by \cref{def:p15-lp-convergence}.
\end{proof}
\end{annexitem}

\begin{annexitem}{annex:p15-w02c-p15-001-chebyshev-square}{thm:p15-w02c-p15-001-chebyshev-square}{The squared form and its centred version}
\begin{proof}
\paragraph{Part (1).}
Let \(a\geq0\). The function \(\varphi(x)=x^2\) is continuous, hence
Borel by \cref{thm:p15-measurable-operations}, and nonnegative. The set
\(A=\{x:|x|\geq a\}=(-\infty,-a]\cup[a,\infty)\) is Borel and contains
\(a\), so it is nonempty. For \(y\in A\), \(|y|\geq a\geq0\) gives
\(y^2=|y|^2\geq a^2\), and \(y=a\) attains the value \(a^2\); hence
\(c_A=\inf_{y\in A}y^2=a^2\). Since \(\{X\in A\}=\{|X|\geq a\}\), the
first display of \cref{thm:p15-expectation-bounds} gives
\[
 a^2\,\mathbb P(|X|\geq a)\leq\mathbb E[X^2;|X|\geq a]\leq\mathbb E[X^2],
\]
where \(\mathbb E[X^2]\in[0,\infty]\) because \(X^2\geq0\).

\paragraph{Part (2).}
Suppose \(\mathbb E[X^2]<\infty\). By \cref{thm:p15-expectation-bounds},
\(X\) is integrable, so \(m\defeq\mathbb E[X]\) is a real number and
\(Y\defeq X-m\) is a real random variable. By
\cref{def:p15-expectation-integral},
\(\mathbb E[Y^2]=\mathbb E[(X-m)^2]=\operatorname{Var}(X)\). Let \(r>0\).
Part~(1) applied to \(Y\) with \(a=r\) gives
\(r^2\,\mathbb P(|X-m|\geq r)\leq\operatorname{Var}(X)\), and dividing by
\(r^2>0\) gives the claim.

\paragraph{Part (3).}
Let \(a<0\) and \(X\equiv0\). Then \(\{|X|\geq a\}=\Omega\), since
\(|X|=0>a\), so the left side is \(a^2\,\mathbb P(\Omega)=a^2>0\). The
right side is \(\mathbb E[0]=0\). For this \(a\) the set
\(A=\{x:|x|\geq a\}\) is all of \(\mathbb R\), and
\(c_A=\inf_{y\in\mathbb R}y^2=0\), attained at \(y=0\), which is why the
argument of part~(1) gives only the true bound \(0\leq\mathbb E[X^2]\).
\end{proof}
\end{annexitem}

\begin{annexitem}{annex:p15-w02c-p15-001-affine-integral}{thm:p15-w02c-p15-001-affine-integral}{The integral of an affine function over an interval}
\begin{proof}
Write \(\iota(x)=x\). It is continuous, hence Borel by
\cref{thm:p15-measurable-operations}, and so is \(\alpha\iota+\beta\).

\paragraph{Integrability.}
Let \(M\defeq\max(|a|,|b|)\). On \([a,b]\),
\(|\alpha x+\beta|\leq|\alpha|M+|\beta|\), so
\(|\alpha\iota+\beta|\mathbf1_{[a,b]}\leq(|\alpha|M+|\beta|)\mathbf1_{[a,b]}\),
and the order statement of \cref{thm:p15-nonnegative-integral-laws} with
part~(3) of \cref{thm:p15-w02c-p15-001-interval-integrals} gives
\(\int_a^b|\alpha x+\beta|\,dx\leq(|\alpha|M+|\beta|)(b-a)<\infty\). In
particular, with \(\alpha=1\) and \(\beta=0\), the function \(\iota\) is
integrable over \([a,b]\).

\paragraph{The case \(a=b\).}
Then \([a,b]=\{a\}\) is \(\lambda\)-null, so
\(\iota\mathbf1_{[a,a]}=0\) a.e., and the a.e.\ replacement statement of
\cref{thm:p15-signed-integral-laws} gives \(\int_a^ax\,dx=0=(a^2-a^2)/2\).

\paragraph{Step functions below and above, when \(a<b\).}
Fix an integer \(n\geq1\), put \(\delta\defeq(b-a)/n>0\) and
\(x_k\defeq a+k\delta\) for \(0\leq k\leq n\), so \(x_0=a\) and
\(x_n=b\). The intervals \(I_k\defeq[x_k,x_{k+1})\),
\(0\leq k\leq n-1\), are disjoint, have union \([a,b)\), and have
\(\lambda(I_k)=\delta\). Define
\[
 s_n\defeq\sum_{k=0}^{n-1}x_k\mathbf1_{I_k},\qquad
 t_n\defeq\sum_{k=0}^{n-1}x_{k+1}\mathbf1_{I_k}.
\]
For \(x\in I_k\), \(s_n(x)=x_k\leq x<x_{k+1}=t_n(x)\); for
\(x\notin[a,b)\), all of \(s_n(x)\), \(\iota(x)\mathbf1_{[a,b)}(x)\) and
\(t_n(x)\) are zero. Hence \(s_n\leq\iota\mathbf1_{[a,b)}\leq t_n\)
everywhere. Each \(\mathbf1_{I_k}\) is integrable with
\(\int\mathbf1_{I_k}\,d\lambda=\lambda(I_k)=\delta\)
(\cref{thm:p15-simple-integral-laws}), so by the linearity and order
statements of \cref{thm:p15-signed-integral-laws},
\[
 \delta\sum_{k=0}^{n-1}x_k
 =\int s_n\,d\lambda
 \leq\int_{[a,b)}\iota\,d\lambda
 \leq\int t_n\,d\lambda
 =\delta\sum_{k=0}^{n-1}x_{k+1}.
\]

\paragraph{Evaluating the two sums.}
For every integer \(n\geq1\), \(\sum_{k=0}^{n-1}k=n(n-1)/2\): this holds
for \(n=1\), and passing from \(n\) to \(n+1\) adds \(n\), with
\(n(n-1)/2+n=(n+1)n/2\). Hence
\[
 \begin{aligned}
 \sum_{k=0}^{n-1}x_k&=na+\delta\frac{n(n-1)}2,\\
 \sum_{k=0}^{n-1}x_{k+1}&=\sum_{k=0}^{n-1}x_k+(x_n-x_0)
 =na+\delta\frac{n(n-1)}2+(b-a).
 \end{aligned}
\]
With \(n\delta=b-a\), the lower and upper integrals are
\[
 a(b-a)+\frac{(b-a)^2}2\cdot\frac{n-1}n
 \qquad\text{and}\qquad
 a(b-a)+\frac{(b-a)^2}2\cdot\frac{n-1}n+\frac{(b-a)^2}n .
\]
Both tend to \(a(b-a)+(b-a)^2/2\) as \(n\to\infty\), and the number
\(\int_{[a,b)}\iota\,d\lambda\) lies between them for every \(n\). Hence
\[
 \int_{[a,b)}\iota\,d\lambda=a(b-a)+\frac{(b-a)^2}2
 =\frac{(b-a)(2a+b-a)}2=\frac{b^2-a^2}2 .
\]
Since \(\iota\) is integrable over \([a,b]\), part~(1) of
\cref{thm:p15-w02c-p15-001-interval-integrals} identifies the left side
with \(\int_a^bx\,dx\).

\paragraph{The affine function.}
The functions \(\iota\mathbf1_{[a,b]}\) and \(\mathbf1_{[a,b]}\) are
integrable, so linearity (\cref{thm:p15-signed-integral-laws}) and
part~(3) of \cref{thm:p15-w02c-p15-001-interval-integrals} give
\[
 \int_a^b(\alpha x+\beta)\,dx
 =\alpha\int\iota\mathbf1_{[a,b]}\,d\lambda+\beta\int\mathbf1_{[a,b]}\,d\lambda
 =\alpha\,\frac{b^2-a^2}2+\beta(b-a).
\]
\end{proof}
\end{annexitem}

\begin{annexitem}{annex:p15-w02c-p15-001-signed-density-values}{thm:p15-w02c-p15-001-signed-density-values}{On the signed density}
\begin{proof}
\paragraph{The three interval integrals.}
By \cref{thm:p15-w02c-p15-001-affine-integral}:
with \(\alpha=2\), \(\beta=-1\) on \([0,1]\),
\(\int_0^1(2x-1)\,dx=2\cdot\frac{1-0}2-1\cdot(1-0)=0\);
with the same \(\alpha,\beta\) on \([1/2,1]\),
\(\int_{1/2}^1(2x-1)\,dx=2\cdot\frac{1-1/4}2-\frac12=\frac34-\frac12=\frac14\);
and with \(\alpha=-2\), \(\beta=1\) on \([0,1/2]\),
\(\int_0^{1/2}(1-2x)\,dx=-2\cdot\frac{1/4-0}2+\frac12=-\frac14+\frac12=\frac14\).

\paragraph{The positive and negative parts.}
For real \(x\), \(h^+(x)=\max(2x-1,0)\) equals \(2x-1\) when
\(x\geq1/2\) and \(0\) otherwise, and \(h^-(x)=\max(1-2x,0)\) equals
\(1-2x\) when \(x\leq1/2\) and \(0\) otherwise. Hence, pointwise,
\[
 h^+\mathbf1_{[0,1]}=h\,\mathbf1_{[1/2,1]},\qquad
 h^-\mathbf1_{[0,1]}=(1-2\iota)\mathbf1_{[0,1/2]},
\]
with \(\iota(x)=x\); at \(x=1/2\) all four functions vanish. The
right-hand functions are nonnegative, and their integrals are
\(\int_{1/2}^1(2x-1)\,dx=1/4\) and \(\int_0^{1/2}(1-2x)\,dx=1/4\) by the
first paragraph and \cref{def:p15-w02c-p15-001-dx-notation}; for
nonnegative functions the signed integral is the nonnegative integral.
So \(\int_{[0,1]}h^+\,d\lambda=\int_{[0,1]}h^-\,d\lambda=1/4\). Since
\(|h|=h^++h^-\), additivity in \cref{thm:p15-nonnegative-integral-laws}
gives \(\int_{[0,1]}|h|\,d\lambda=1/2\). As a check,
\cref{def:p15-integrable-functions} gives
\(\int_{[0,1]}h\,d\lambda=1/4-1/4=0\), agreeing with the first integral.
\end{proof}
\end{annexitem}

\begin{annexitem}{annex:p15-w02c-p15-001-reciprocal-density}{thm:p15-w02c-p15-001-reciprocal-density}{Reciprocal}
\begin{proof}
\paragraph{Measurability.}
On \((0,1)\), \(q(x)=1/x>1\); elsewhere \(q(x)=0\). For real \(t\):
if \(t<0\), \(\{q\leq t\}=\varnothing\); if \(0\leq t\leq1\),
\(\{q\leq t\}=(-\infty,0]\cup[1,\infty)\); if \(t>1\), then for
\(0<x<1\) the inequality \(1/x\leq t\) holds exactly when \(x\geq1/t\),
and \(1/t<1\), so \(\{q\leq t\}=(-\infty,0]\cup[1/t,\infty)\). Each of
these sets is closed, hence Borel. Since the rays \((-\infty,t]\) generate
\(\mathcal B(\mathbb R)\) (\cref{thm:p15-borel-generators}),
\cref{thm:p15-measurability-criterion} shows that \(q\) is Borel.

\paragraph{Part (1).}
Let \(n\geq2\). For \(x\in[1/n,1)\), \(q(x)=1/x\leq n\). Hence
\(q\mathbf1_{[1/n,1)}\leq n\mathbf1_{[1/n,1)}\), and the order statement
of \cref{thm:p15-nonnegative-integral-laws} with
\cref{thm:p15-simple-integral-laws} gives
\[
 \int_{[1/n,1)}q\,d\lambda\leq n\,\lambda([1/n,1))=n\left(1-\frac1n\right)=n-1 .
\]

\paragraph{Part (2).}
Let \(0<a\leq1\). For integers \(k\geq0\) put
\(J_k\defeq[a2^{-k-1},a2^{-k})\). These intervals are pairwise disjoint,
since for \(j<k\) every point of \(J_k\) is less than
\(a2^{-k}\leq a2^{-j-1}\), the left endpoint of \(J_j\). Their union is
\((0,a)\): each \(J_k\) lies in \((0,a)\); and for \(x\in(0,a)\), since
\(a2^{-k-1}\to0\), there is a least integer \(k\geq0\) with
\(a2^{-k-1}\leq x\), and then \(x<a2^{-k}\), by \(x<a\) when \(k=0\) and
by the minimality of \(k\) when \(k\geq1\). So \(x\in J_k\).

For \(x\in J_k\), \(0<x<a2^{-k}\leq1\), so \(q(x)=1/x>2^k/a\). Hence
\(q\mathbf1_{J_k}\geq(2^k/a)\mathbf1_{J_k}\), and
\[
 \int_{J_k}q\,d\lambda\geq\frac{2^k}a\,\lambda(J_k)
 =\frac{2^k}a\cdot a\,2^{-k-1}=\frac12 .
\]
Part~(4) of \cref{thm:p15-nonnegative-integral-laws} makes
\(E\mapsto\int_Eq\,d\lambda\) a measure, so
\[
 \int_{(0,a)}q\,d\lambda=\sum_{k\geq0}\int_{J_k}q\,d\lambda
 \geq\sum_{k\geq0}\frac12=\infty .
\]
For the last claim, take \(a=1\): \([0,1]\) and \((0,1)\) differ by the
\(\lambda\)-null set \(\{0,1\}\), on which \(q\) vanishes anyway, so
\(\int_0^1q(x)\,dx=\int_{(0,1)}q\,d\lambda=\infty\) by part~(1) of
\cref{thm:p15-w02c-p15-001-interval-integrals}, applied in its
nonnegative case.
\end{proof}
\end{annexitem}

\begin{annexitem}{annex:p15-w02c-p15-001-lower-moments}{thm:p15-w02c-p15-001-lower-moments}{Lower absolute moments}
\begin{proof}
\paragraph{Measurability.}
For real \(r>0\), the map \(t\mapsto t^r\) on \([0,\infty)\) is
continuous: on \((0,\infty)\) it is \(t\mapsto\exp(r\log t)\), a
composite of continuous maps, and at \(0\) it is continuous because
\(0\leq t^r\leq\varepsilon\) whenever \(0\leq t\leq\varepsilon^{1/r}\).
Composing with the continuous map \(x\mapsto|x|\) shows that
\(x\mapsto|x|^r\) is continuous on \(\mathbb R\), hence Borel
(\cref{thm:p15-measurable-operations}). So \(|X|^j\) and \(|X|^k\) are
nonnegative random variables, and their expectations exist in
\([0,\infty]\) (\cref{def:p15-expectation-integral}).

\paragraph{The bound.}
Put \(p\defeq k/j\) and \(q\defeq k/(k-j)\). Since \(0<j<k\), both lie
in \((1,\infty)\), and \(1/p+1/q=j/k+(k-j)/k=1\). Apply
\cref{thm:p15-holder-integral} on \((\Omega,\mathcal F,\mathbb P)\) to
the real measurable functions \(f\defeq|X|^j\) and \(g\defeq1\). For
\(t\geq0\), \((t^j)^{k/j}=t^k\), so \(|f|^p=|X|^k\) and
\[
 \|f\|_p=\left(\mathbb E|X|^k\right)^{1/p}
 =\left(\mathbb E|X|^k\right)^{j/k},
\]
with the value \(\infty\) when \(\mathbb E|X|^k=\infty\). Also
\(|g|^q=1\), so \(\|g\|_q=\mathbb P(\Omega)^{1/q}=1\). Since
\(|fg|=|X|^j\), H\"older's inequality reads
\[
 \mathbb E|X|^j=\int|fg|\,d\mathbb P\leq\|f\|_p\|g\|_q
 =\left(\mathbb E|X|^k\right)^{j/k}.
\]
If \(\mathbb E|X|^k<\infty\), the right side is a real number, so
\(\mathbb E|X|^j<\infty\).

\paragraph{The norm form.}
If \(\mathbb E|X|^k=\infty\), the right side of the second display of the
theorem is \(\infty\) and there is nothing to prove. Otherwise both
\(\mathbb E|X|^j\) and \((\mathbb E|X|^k)^{j/k}\) are nonnegative real
numbers, and \(t\mapsto t^{1/j}\) is nondecreasing on \([0,\infty)\).
Applying it to the bound gives
\((\mathbb E|X|^j)^{1/j}\leq((\mathbb E|X|^k)^{j/k})^{1/j}
=(\mathbb E|X|^k)^{1/k}\). Conversely, applying the nondecreasing map
\(t\mapsto t^j\) to the norm form returns the first bound, so the two forms
are equivalent. For \(1\leq j<k\) the two sides are \(\|X\|_j\) and
\(\|X\|_k\) by \cref{def:p15-lp-convergence}.
\end{proof}
\end{annexitem}

\begin{annexitem}{annex:p15-w02c-p15-001-lower-moments-limits}{thm:p15-w02c-p15-001-lower-moments-limits}{Limits of the moment bound}
\begin{proof}
\paragraph{Part (1).}
Every function on \(I\) is measurable for the sigma-algebra of all subsets.
Use the enumeration \(i_k=k\) in part~(2) of
\cref{thm:p15-w02c-p15-001-counting-sums}. For \(i\geq2\),
\(i^2\geq i(i-1)\), so
\[
 \frac1{i^2}\leq\frac1{i(i-1)}=\frac1{i-1}-\frac1i .
\]
For \(N\geq2\) the sum of the right sides over \(2\leq i\leq N\)
telescopes to \(1-1/N\), so
\(\sum_{i\leq N}1/i^2\leq1+1-1/N\leq2\); the same bound holds for
\(N=1\). Hence \(\int f^2\,d\#=\lim_N\sum_{i\leq N}1/i^2\leq2\).
For the first power, group the indices into the blocks
\(B_m=\{2^m+1,\ldots,2^{m+1}\}\), \(m\geq0\). The block \(B_m\) has
\(2^m\) members, each at most \(2^{m+1}\), so
\(\sum_{i\in B_m}1/i\geq2^m\cdot2^{-m-1}=1/2\). For \(M\geq1\), the
indices \(1,\ldots,2^M\) are \(1\) together with \(B_0,\ldots,B_{M-1}\),
so \(\sum_{i\leq2^M}1/i\geq1+M/2\). The partial sums are nondecreasing
and unbounded, so \(\int f\,d\#=\lim_N\sum_{i\leq N}1/i=\infty\).

\paragraph{Part (2).}
The masses \(p(n)=2^{-n}\) on \(I=\{1,2,\ldots\}\) have sum one by the
finite geometric identity of \cref{def:p15-countable-waiting}, so
\(\mathbb P=\mu_p\) is a probability measure on all subsets of \(I\) by
\cref{thm:p15-countable-atomic}. The variable \(X(n)=2^n\) is measurable,
and so are \(X^{1/2}\) and \(X\), which are nonnegative. On this countable
atomic space, \cref{thm:p15-expectation-law-integral} evaluates a
nonnegative expectation as \(\sum_nY(n)p(n)\). Hence
\(\mathbb EX=\sum_{n\geq1}2^n2^{-n}=\sum_{n\geq1}1=\infty\). Next let
\(r\defeq2^{-1/2}\in(0,1)\); then
\(X(n)^{1/2}p(n)=2^{n/2}2^{-n}=r^n\). For \(N\geq1\),
\((1-r)\sum_{n=1}^Nr^n=r-r^{N+1}\), so
\(\sum_{n=1}^Nr^n=r(1-r^N)/(1-r)\). Since \(0<r<1\), \(r^N\to0\), and the
partial sums increase to \(r/(1-r)\). Therefore
\[
 \mathbb EX^{1/2}=\frac r{1-r}=\frac{2^{-1/2}}{1-2^{-1/2}}
 =\frac1{\sqrt2-1}<\infty .
\]

\paragraph{Part (3).}
Let \(N\) be the null event \(\{|X|\neq c\}\). Off \(N\),
\(|X|^j=c^j\) and \(|X|^k=c^k\). By part~(2) of
\cref{thm:p15-nonnegative-integral-laws} and the integral of a constant
nonnegative simple function (\cref{def:p15-simple-integral}), \(\mathbb E|X|^j=c^j\,\mathbb P(\Omega)=c^j\)
and \(\mathbb E|X|^k=c^k\). Finally \((c^k)^{j/k}=c^j\) for \(c\geq0\).
\end{proof}
\end{annexitem}

\begin{annexitem}{annex:p15-w02c-p15-001-bernoulli-moments}{thm:p15-w02c-p15-001-bernoulli-moments}{Bernoulli moments}
\begin{proof}
\paragraph{A null exceptional event.}
The singletons \(\{0\}\) and \(\{1\}\) are Borel, so
\(A\defeq\{X=1\}\) and \(\{X=0\}\) are events. They are disjoint, and
\(\mathbb P(A\cup\{X=0\})=p+(1-p)=1\). Hence
\(N\defeq\{X\notin\{0,1\}\}\), the complement of their union, has
\(\mathbb P(N)=0\).

\paragraph{Pointwise identities off \(N\).}
Let \(\omega\notin N\). If \(X(\omega)=1\), then \(\omega\in A\) and
\(\mathbf1_A(\omega)=1\); if \(X(\omega)=0\), then \(\omega\notin A\)
and \(\mathbf1_A(\omega)=0\). So \(X=\mathbf1_A\) off \(N\). Consequently,
off \(N\), \(X^+=\mathbf1_A\), \(X^-=0\) and \(X^2=\mathbf1_A=X\), since
\(t^2=t\) for \(t\in\{0,1\}\). Thus \(X^2=X\) almost surely.

\paragraph{The moments.}
By part~(2) of \cref{thm:p15-nonnegative-integral-laws}, nonnegative
functions that agree a.e.\ have equal integrals. Hence
\(\mathbb EX^+=\mathbb E\mathbf1_A=\mathbb P(A)=p\) and
\(\mathbb EX^-=0\). Both are finite, so \(X\) is integrable
(\cref{def:p15-integrable-functions}) and
\(\mathbb EX=\mathbb EX^+-\mathbb EX^-=p\). The nonnegative variable
\(X^2\) agrees a.e.\ with \(\mathbf1_A\), so \(\mathbb EX^2=p<\infty\).

\paragraph{The variance.}
Since \(\mathbb EX^2<\infty\), the variance identity of
\cref{thm:p15-expectation-bounds} applies:
\(\operatorname{Var}(X)=\mathbb EX^2-(\mathbb EX)^2=p-p^2=p(1-p)\).
Expanding, \(\tfrac14-(p-\tfrac12)^2=\tfrac14-p^2+p-\tfrac14=p-p^2\).
As \((p-\tfrac12)^2\geq0\), with equality exactly when \(p=1/2\),
\(p(1-p)\leq\tfrac14\) with equality exactly then. Finally
\(p(1-p)=0\) exactly when \(p=0\) or \(p=1\).
\end{proof}
\end{annexitem}

\begin{annexitem}{annex:p15-w02c-p15-001-jensen-norm}{thm:p15-w02c-p15-001-jensen-normalization}{Total mass one}
\begin{proof}
\paragraph{(1) implies (2).}
If \(\mu(\Omega)=1\), then \(\mu\) is a probability measure, and the
inequality in~(2) is \cref{thm:p15-jensen-integral}.

\paragraph{The constant test functions.}
For a real constant \(k\), the function \(\varphi\equiv k\) is convex,
because \(k\leq tk+(1-t)k\) holds with equality for \(0<t<1\). The zero
function \(f\equiv0\) is a nonnegative simple function with integral
\(0\), so it is integrable with \(\int f\,d\mu=0\)
(\cref{def:p15-simple-integral,def:p15-integrable-functions}). For
\(k=1\), \(\varphi(f)\) is the constant simple function \(1\), whose
integral is \(\mu(\Omega)<\infty\). For \(k=-1\), \(\varphi(f)\equiv-1\)
has positive part \(0\) and negative part \(1\), so it is integrable and
\(\int\varphi(f)\,d\mu=0-\mu(\Omega)=-\mu(\Omega)\) by
\cref{def:p15-integrable-functions}.

\paragraph{(2) implies (1).}
Assume~(2). Applied to \(f\equiv0\) and \(\varphi\equiv1\), it gives
\(\varphi(0)=1\leq\mu(\Omega)\). Applied to \(f\equiv0\) and
\(\varphi\equiv-1\), it gives \(-1\leq-\mu(\Omega)\), that is,
\(\mu(\Omega)\leq1\). Hence \(\mu(\Omega)=1\).

\paragraph{The violations.}
The same two computations show that, if \(\mu(\Omega)<1\), the pair
\((0,1)\) gives \(\varphi(\int f\,d\mu)=1>\mu(\Omega)=\int\varphi(f)\,d\mu\),
and that, if \(1<\mu(\Omega)<\infty\), the pair \((0,-1)\) gives
\(\varphi(\int f\,d\mu)=-1>-\mu(\Omega)=\int\varphi(f)\,d\mu\).
\end{proof}
\end{annexitem}

\begin{annexitem}{annex:p15-w02c-p15-001-jensen-infinite-mass}{thm:p15-w02c-p15-001-jensen-infinite-mass}{Jensen: infinite mass}
\begin{proof}
\paragraph{Part (1).}
The function \(x\mapsto x^2\) is convex: for real \(x,y\) and
\(0<t<1\),
\[
 tx^2+(1-t)y^2-(tx+(1-t)y)^2=t(1-t)(x-y)^2\geq0,
\]
as expanding both sides shows. The indicator \(f=\mathbf1_A\) is a
nonnegative simple function with \(\int f\,d\mu=\mu(A)<\infty\), so it is
integrable. Since \(1^2=1\) and \(0^2=0\), \(\varphi(f)=f\), which is
integrable with the same integral \(\mu(A)\). Then
\(\varphi(\int f\,d\mu)=\mu(A)^2\), and \(\mu(A)^2>\mu(A)\) because
\(\mu(A)>1\). On \((\mathbb R,\mathcal B(\mathbb R),\lambda)\),
\(\lambda((0,2])=2\) by \cref{thm:p15-lebesgue-unit-interval}, and
\(\lambda(\mathbb R)=\infty\) by the same theorem, so \(A=(0,2]\)
qualifies and gives \(2^2=4>2\).

\paragraph{Part (2): integrable functions vanish a.e.}
Let \(f\) be real and integrable, and let \(n\geq1\). Put
\(B_n\defeq\{|f|>1/n\}\in\mathcal F\). Pointwise
\(\mathbf1_{B_n}\leq n|f|\), so by monotonicity of the nonnegative
integral (\cref{thm:p15-monotone-integral}) and scaling
(\cref{thm:p15-nonnegative-integral-laws}),
\(\mu(B_n)\leq n\int|f|\,d\mu<\infty\). By hypothesis \(\mu(B_n)\) is
\(0\) or \(\infty\), so \(\mu(B_n)=0\). The set \(\{f\neq0\}\) is the
union of the \(B_n\), so it is null by countable subadditivity. Thus
\(f=0\) a.e., and \(\int f\,d\mu=0\) by
\cref{thm:p15-signed-integral-laws}, since a.e.\ replacement by the zero
function preserves the integral.

\paragraph{Part (2): the admissible \(\varphi\).}
Let \(Z\defeq\{f=0\}\). Since \(\Omega=Z\cup\{f\neq0\}\) and the second
set is null, \(\mu(Z)=\mu(\Omega)=\infty\). On \(Z\),
\(|\varphi(f)|=|\varphi(0)|\), so
\(|\varphi(f)|\geq|\varphi(0)|\mathbf1_Z\) and
\(\int|\varphi(f)|\,d\mu\geq|\varphi(0)|\,\mu(Z)\). If
\(\varphi(0)\neq0\) the right side is \(\infty\) and \(\varphi(f)\) is not
integrable. So an admissible pair has \(\varphi(0)=0\); then
\(\varphi(f)=\varphi(0)=0\) a.e., \(\int\varphi(f)\,d\mu=0\), and
\(\varphi(\int f\,d\mu)=\varphi(0)=0\). The inequality holds with
equality.

\paragraph{The one-point instance.}
On \(\{\omega_0\}\) the only measurable sets are \(\varnothing\) and the
whole space, with measures \(0\) and \(\infty\), and the whole space has
infinite measure, so part~(2) applies.
\end{proof}
\end{annexitem}

\begin{annexitem}{annex:p15-w02c-p15-001-jensen-extended}{thm:p15-w02c-p15-001-jensen-extended}{Jensen: extended form}
\begin{proof}
\paragraph{Measurability and the supporting function.}
The proof of \cref{thm:p15-jensen-integral} shows that a convex
\(\varphi\) is continuous, hence Borel measurable, so \(\varphi(X)\) is a
random variable. The same proof constructs, for each real \(c\), a real
\(a\) with
\[
 \varphi(x)\geq\ell(x)\defeq\varphi(c)+a(x-c)\qquad\text{for every real }x.
\]
Take \(c\defeq\mathbb EX\), which is real because \(X\) is integrable.

\paragraph{The negative part is integrable.}
Pointwise \(\varphi(X)\geq\ell(X)\), so \(-\varphi(X)\leq-\ell(X)\) and
\(\varphi(X)^-=\max(-\varphi(X),0)\leq\max(-\ell(X),0)\leq|\ell(X)|\).
By the triangle inequality,
\(|\ell(X)|\leq|\varphi(c)|+|a||c|+|a||X|\). The right side is a
nonnegative random variable with expectation
\(|\varphi(c)|+|a||c|+|a|\,\mathbb E|X|<\infty\), by additivity and
scaling of nonnegative expectations
(\cref{thm:p15-nonnegative-integral-laws}). Monotonicity
(\cref{thm:p15-monotone-integral}) gives
\(\mathbb E[\varphi(X)^-]<\infty\). By
\cref{def:p15-integrable-functions}, \(\mathbb E\varphi(X)\) exists and
equals \(\mathbb E[\varphi(X)^+]-\mathbb E[\varphi(X)^-]\in(-\infty,\infty]\).

\paragraph{The inequality.}
If \(\mathbb E[\varphi(X)^+]=\infty\), then
\(\mathbb E\varphi(X)=\infty\geq\varphi(c)\). Otherwise both parts have
finite expectation, \(\varphi(X)\) is integrable, and
\cref{thm:p15-jensen-integral} with \(\mu=\mathbb P\) and \(f=X\) gives
\(\varphi(\mathbb EX)\leq\mathbb E\varphi(X)\). The two cases are
exhaustive, and in the first \(\varphi(X)\) is not integrable.
\end{proof}
\end{annexitem}

\begin{annexitem}{annex:p15-w02c-p15-001-jensen-integrability-limits}{thm:p15-w02c-p15-001-jensen-integrability-limits}{Jensen: integrability}
\begin{proof}
\paragraph{Sums against masses.}
Let \(\Omega\) be countable with all subsets measurable, let
\(r:\Omega\to[0,\infty)\) have sum one, and let \(\mathbb P=\mu_r\) be the
probability measure of \cref{thm:p15-countable-atomic}. For nonnegative
\(Y\) on \(\Omega\), enumerate \(\Omega\) as \(\omega_1,\omega_2,\ldots\);
then \(Y=\sum_kY(\omega_k)\mathbf1_{\{\omega_k\}}\) pointwise, a series of
nonnegative simple functions with integrals \(Y(\omega_k)r(\omega_k)\).
Term-by-term integration of nonnegative series
(\cref{thm:p15-nonnegative-integral-laws}) gives
\(\mathbb EY=\sum_kY(\omega_k)r(\omega_k)\) in \([0,\infty]\), and by
the grouping part of \cref{thm:p15-countable-atomic} the value does not
depend on the enumeration and may be computed in groups.

\paragraph{Part (1).}
Let \(S_N\defeq\sum_{n=1}^N4^{-n}\). Then
\(S_N-\tfrac14S_N=\tfrac14-4^{-N-1}\), so \(S_N=(1-4^{-N})/3\), and the
partial sums of \(\sum_n3\cdot4^{-n}\) increase to \(1\). Hence
\(\mathbb P\) is a probability measure on \(I\). By the first paragraph,
\[
 \mathbb EX=\sum_{n\geq1}2^n\cdot3\cdot4^{-n}=3\sum_{n\geq1}2^{-n}=3,
\]
using \(\sum_{n\geq1}2^{-n}=1\) from \cref{def:p15-countable-waiting},
so \(X\geq0\) is integrable with \(\mathbb EX=3\); and
\[
 \mathbb E(X^2)=\sum_{n\geq1}4^n\cdot3\cdot4^{-n}=\sum_{n\geq1}3=\infty .
\]
So \(\varphi(X)=X^2\) is not integrable,
\cref{thm:p15-jensen-integral} does not apply, and
\cref{thm:p15-w02c-p15-001-jensen-extended} applies with
\(\varphi(\mathbb EX)=9\) and \(\mathbb E\varphi(X)=\infty\).

\paragraph{Part (2).}
Group \(J\) into the pairs \(\{n,-n\}\), \(n\geq1\). Each pair has mass
\(2\cdot2^{-n-1}=2^{-n}\), so by the grouping part of
\cref{thm:p15-countable-atomic} the total mass is
\(\sum_{n\geq1}2^{-n}=1\), and \(\mathbb P\) is a probability measure. The
positive part is \(X^+(m)=2^m\) for \(m>0\) and \(0\) for \(m<0\); the
negative part is \(X^-(m)=2^{-m}\) for \(m<0\) and \(0\) for \(m>0\). By
the first paragraph,
\[
 \mathbb E[X^+]=\sum_{m\geq1}2^m2^{-m-1}=\sum_{m\geq1}\tfrac12=\infty,
 \qquad
 \mathbb E[X^-]=\sum_{n\geq1}2^n2^{-n-1}=\infty .
\]
Both parts have infinite expectation, so by
\cref{def:p15-integrable-functions} the signed integral of \(X\) is not
defined, and there is no real number at which to evaluate \(\varphi\).
\end{proof}
\end{annexitem}
